\section{Runtime loop and modules}
\label{sec:runtime}

\subsection{CellMain}
\pathref{App/CellMain.cpp} owns logger init/shutdown and the frame loop:
\begin{lstlisting}
// conceptual
Logger::initLogger();
Illumo* illumo = new Illumo(argc, argv);
illumo->Init();
while (!illumo->ShouldClose()) {
  double dt = /* glfw time delta */;
  illumo->Update(dt);
  illumo->Render();
}
illumo->Shutdown();
\end{lstlisting}

\subsection{Illumo as composition root}
\symbolref{Illumo} owns the long-lived services (\symbolref{EnvVars},
\symbolref{RenderWindow}, \symbolref{Camera}, \symbolref{Renderer},
\symbolref{AssetManager}, \symbolref{CommandRegistry}, \symbolref{CommandLine},
\symbolref{InputManager}, \symbolref{Scene}) and a vector of
\symbolref{IModule} instances.

\symbolref{IllumoContext} is a \textbf{raw-pointer bag} into those services,
passed into modules at \symbolref{Start}. Ownership stays on \symbolref{Illumo}.

\begin{inferencebox}
Pattern: composition root + service locator-ish context struct (not a full
service locator library). Modules are closer to subsystems than ECS systems.
\end{inferencebox}

\subsection{IModule contract}
\begin{lstlisting}
class IModule {
public:
  virtual void Start(IllumoContext* context) = 0;
  virtual void Update(double dt) = 0;
  virtual void DispatchDrawables(Scene* scene) = 0;
  virtual void Exit() = 0;
};
\end{lstlisting}

\begin{itemize}
  \item \textbf{Start} --- bind inputs, allocate game state.
  \item \textbf{Update} --- simulation / input.
  \item \textbf{DispatchDrawables} --- contribute what should draw this frame.
  \item \textbf{Exit} --- teardown module-owned state.
\end{itemize}

Registered modules today (from \symbolref{Illumo::Init}):
\begin{itemize}
  \item \symbolref{CellGameModule} (always)
  \item \symbolref{DebugModule} (non-\texttt{NDEBUG})
\end{itemize}

\begin{openbox}
Engine header \pathref{Engine/Illumo.h} currently includes concrete
\pathref{Game/CellGameModule.h}. Intended long-term: register modules from App
or a bootstrap file so Engine does not depend on Game.
\end{openbox}

\subsection{Frame flow (conceptual)}
\begin{center}
\begin{tikzpicture}[
  font=\scriptsize,
  node distance=0.55cm,
  % Do not name this style ``step'' --- it clashes with TikZ's /tikz/step key.
  flowstep/.style={draw, rounded corners, align=center, fill=blue!6, minimum width=4.2cm},
  arr/.style={-{Latex}}
]
  \node[flowstep] (1) {InputManager::update};
  \node[flowstep, below=of 1] (2) {Camera::Update};
  \node[flowstep, below=of 2] (3) {foreach module: Update(dt)};
  \node[flowstep, below=of 3] (4) {Scene clear + modules DispatchDrawables};
  \node[flowstep, below=of 4] (5) {Renderer::RenderScene (tokens + submit)};
  \node[flowstep, below=of 5] (6) {IBackend::EndFrame (swap)};
  \foreach \a/\b in {1/2,2/3,3/4,4/5,5/6} \draw[arr] (\a) -- (\b);
\end{tikzpicture}
\end{center}
